\clearpage
\fancyhead[L]{\small EIN NEUES HILFSMITTEL.}
\SourceLeaf{358}
\pdfbookmark[0]{Göttingische gelehrte Anzeigen. 30. Oktober 1837.}{anzeige-1837-10-30}
\begin{center}
Göttingische gelehrte Anzeigen. 30. Oktober 1837.
\end{center}

Der Vortrag, den Hofrat \textsc{Gauss} am 19. September in der Sitzung der Königlichen Sozietät hielt, behandelt
\begin{center}
\textit{ein neues Hilfsmittel für die magnetischen Beobachtungen.}
\end{center}

Die magnetische Deklination ist eine der Größen, in denen sich der Erdmagnetismus äußert. Sie hat die Beobachter am frühesten beschäftigt und wird den anderen Größen in mehrfacher Hinsicht weiterhin vorgezogen, seit sich die Aufmerksamkeit der Naturforscher auch diesen zugewandt hat. Eine der interessantesten Untersuchungen auf dem Gebiet des Erdmagnetismus betrifft die erstaunlichen unregelmäßigen Störungen dieser Kraft, die dennoch über einen ganzen Erdteil gleichzeitig und gleichmäßig wirken. Dafür hat sich in den letzten Jahren ein eigener Verein von Beobachtern zusammengeschlossen, und die beiden anderen Größen sind bisher ganz von dieser Untersuchung ausgeschlossen. Der Vorzug der Deklination erklärt sich weniger aus der vielfältigen praktischen Bedeutung ihrer Kenntnis für Seefahrer, Geodäten und Markscheider als aus dem bisherigen Stand der Beobachtungsmittel. Diese lassen für die Deklination nichts zu wünschen übrig, sind aber für die anderen Größen noch viel weniger entwickelt. Das vor einigen Jahren eingeführte Magnetometer dient zwar neben der Bestimmung der Deklination auch zur Ermitt\SourceLeaf{359}lung der horizontalen Intensität. Gerade die Aufgabe, diese auf ein absolutes Maß zurückzuführen, gab den ersten Anlass zu diesem Instrument. Es löst die Aufgabe jedoch noch nicht in \textit{jeder} Hinsicht. Seiner Natur nach kann es nur einen Mittelwert der Intensität über einen gewissen Zeitraum genau angeben. Obwohl dieser Zeitraum. in gewissem Maß nach Belieben gewählt werden kann, darf man ihn nicht zu kurz ansetzen, weil sich sonst mit ihm auch die Genauigkeit und Zuverlässigkeit des Ergebnisses ändern würden. Das Magnetometer eignet sich daher überhaupt nicht dazu, die Intensitätsänderungen in kurzen Zeitabständen zu verfolgen. Der neue Apparat, von dem der Vortrag handelt, soll diese Lücke für die Intensität schließen und die Beobachtung dieser Größe ebenso leicht und zuverlässig machen, wie es das Magnetometer für die Deklination tut.

Der begrenzte Platz in diesen Blättern erlaubt es nicht, den Apparat zu beschreiben, den Hofrat \textsc{Gauss} hat anfertigen lassen. Das ist umso weniger nötig, als der Vortrag selbst bald im Druck erscheinen wird. Das Instrument ist seit einigen Monaten in der hiesigen Sternwarte aufgehängt. Schon bei zwei magnetischen Beobachtungsterminen wurden daran die Änderungen der Intensität jeweils 24 Stunden lang beobachtet. Trägt man die gleichzeitig im magnetischen Observatorium beobachteten Änderungen der Deklination zusammen mit ihnen in eine einzige Zeichnung ein, so tritt das wunderbare Spiel der magnetischen Störungen auf eine neue, eigentümliche Weise sehr anschaulich hervor. Man kann mit Zuversicht erwarten, dass wir bald umfassendere Aufschlüsse darüber gewinnen werden, wo die Ursachen dieser rätselhaften Erscheinungen liegen, sobald an mehreren weit voneinander entfernten Orten auf ähnliche Weise gleichzeitig beobachtet wird. Die Intensität zeigt ebenso häufige und ebenso beträchtliche unregelmäßige Störungen wie die Deklination. Bei beiden treten aber auch regelmäßige Änderungen hervor, die mit der Tageszeit zusammenhängen. Bei der Intensität verlaufen sie etwas anders, soweit sich dies aus täglichen Aufzeichnungen zu bestimmten Stunden während eines Monats erkennen lässt. Die westliche Deklination nimmt nämlich in unseren Gegenden von 7 oder 8 Uhr morgens bis eine oder zwei Stunden nach Mittag zu und geht dann wieder zurück. Die Intensität nimmt dagegen in den ersten Vormittagsstunden ab, erreicht ihr Minimum aber schon eine oder zwei Stunden vor Mittag, wenn die Deklination ge\SourceLeaf{360}rade am schnellsten zunimmt. Es braucht kaum daran erinnert zu werden, dass an einzelnen Tagen die unregelmäßigen Störungen diesen regelmäßigen Verlauf oft völlig verdecken und dass genauere Bestimmungen erst aus lange fortgesetzten Beobachtungen hervorgehen können.

Die Bauart erlaubt es, den Apparat außer für seinen Hauptzweck noch für viele ganz verschiedene Zwecke zu verwenden. Er löst ein Problem, mit dem man sich früher wiederholt, aber ohne Erfolg beschäftigt hat, nämlich die stündlichen und die unregelmäßigen Änderungen der Deklination vergrößert darzustellen. In der gegenwärtigen Anordnung des Apparats beträgt die Vergrößerung das Zehnfache. Eine Änderung der Deklination, die sich am Magnetometer des magnetischen Observatoriums in 50 Skalenteilen zeigt, erscheint hier mit 500 Skalenteilen. Beim letzten magnetischen Beobachtungstermin (30. September) hat sich dies durch 8 Stunden gleichzeitiger Beobachtung an beiden Apparaten aufs Befriedigendste bestätigt.

Bekanntlich ergibt das Magnetometer zusammen mit einem Multiplikator ein sehr empfindliches Galvanometer. Das gilt auch für den neuen Apparat, dessen Empfindlichkeit hierbei aber die des Magnetometers genau in demselben Verhältnis übertrifft, das wir für Deklinationsänderungen angegeben haben. Der neue Apparat dient also zur genauen Messung selbst der schwächsten galvanischen Ströme. Es erregt gewöhnlich Bewunderung, dass diese den 25-pfündigen Magnetstab in dem Apparat in so beträchtliche Bewegungen versetzen. Bei thermogalvanisch erzeugten Strömen widerlegt dies aufs Deutlichste die irrige Meinung vieler Physiker, dass solche Ströme eine Kette von beträchtlicher Länge nicht durchlaufen könnten. Eine noch so lange Kette hebt solche Ströme nicht auf. Sie schwächt sie nur, und zwar genau im selben Verhältnis wie bei anderen Arten der Erregung. Bei einem thermogalvanischen Apparat besonderer Bauart erzeugt schon die Berührung der Verbindungsstellen mit dem Finger einen galvanischen Strom. Selbst wenn dieser eine fast zwei Meilen lange Kette aus größtenteils sehr dünnem Draht durchlaufen muss, bleibt er an sehr beträchtlichen Ausschlägen des Magnetstabs erkennbar.

Lässt man gewöhnliche Reibungselektrizität durch den Multiplikator gehen, so gehören ihre elektromagnetischen Wirkungen zu den schwächsten. Sie sind schwer zu erkennen und noch schwerer zu messen. Bekanntlich wurde das
\SourceLeaf{361}Vorhandensein solcher Wirkungen zuerst von \textsc{Colladon} entdeckt und später. von \textsc{Faraday} bestätigt. Statt wie diese Physiker eine starke elektrische Batterie durch den Leitungsdraht zu entladen, beobachtete Hofrat \textsc{Gauss} die Wirkung der Reibungselektrizität, während eine im physikalischen Kabinett aufgestellte Elektrisiermaschine fortwährend gedreht wurde. Ihr Konduktor und ihr Reibzeug waren mit den Enden der großen, zur Sternwarte führenden Kette verbunden. In dieser Kette lag der Multiplikator, der den Magnetstab des neuen Apparats umgibt. Der Stab wurde dadurch bei einem Ausschlag von 144 Skalenteilen oder 51 Bogenminuten gehalten, positiv oder negativ, je nachdem, in welcher Richtung die Elektrizität den Multiplikator durchlief. Die Drahtlänge der Kette betrug dabei etwa 13000 Fuß. Besonders bemerkenswert ist, dass eine Verlängerung der Kette bis auf fast eine ganze Meile durch Einfügen weiteren Drahts die elektromagnetische Wirkung überhaupt nicht verringerte. In dieser Hinsicht verhält sich die strömende Maschinenelektrizität anders als die galvanischen Ströme, die hydrogalvanisch, thermogalvanisch oder durch Induktion erzeugt werden. Deren durch die magnetische Wirkung gemessene Intensität wird umso schwächer, je länger die schließende Kette ist. Diese Erscheinung beweist keineswegs einen wesentlichen inneren Unterschied zwischen diesen Strömen. Sie bestätigt vielmehr ihre Gleichartigkeit und die Theorie, nach der ungleiche Intensität zweier Ströme nichts anderes bedeutet als eine ungleiche Menge an Elektrizität, die in gegebener Zeit durch jeden Querschnitt der Leitung strömt. Gegebene elektromotorische Kräfte der zuletzt genannten Arten setzen lediglich umso weniger Elektrizität in Bewegung, je größer der Widerstand ist, den die längere Kette entgegensetzt. Bei dem oben angeführten Versuch musste aber \textit{alle} Elektrizität, die von der Maschine in Form von Funken auf den Konduktor übersprang, die ganze Kette durchlaufen, um sich mit der entgegengesetzten Elektrizität des Reibzeugs auszugleichen. Das galt bei kurzer wie bei langer Kette, sofern sie nur hinreichend isoliert war. Die Menge der Elektrizität, die in bestimmter Zeit jeden Querschnitt des Leitungsdrahts durchströmte, hing also überhaupt nicht von der Länge der Kette ab, sondern nur vom Betrieb der Maschine.

Bei den meisten der hier erwähnten Versuche musste der galvanische oder elektrische Strom die große Kette durchlaufen, die im Jahr 1833 zwischen der Sternwarte und dem physikalischen Ka%
\SourceLeaf{362}binett gezogen wurde. Allein ihr in der Luft verlaufender Draht hat eine Länge von fast 7000 Fuß. Der Hauptzweck dieser Anlage besteht zwar darin, die Gesetze der galvanischen Ströme in großem Maßstab physikalisch zu untersuchen. Doch wurde die Gelegenheit von Anfang an auch vielfach zu Versuchen mit elektromagnetischer Telegrafie genutzt, die selbst mit ganzen Wörtern und kurzen Wendungen aufs Befriedigendste gelangen. Das wurde in diesen Blättern schon bei der ersten Nachricht von der Einrichtung des hiesigen magnetischen Observatoriums erwähnt (Gött. gel. Anz. 1834, 9. August). Das zunächst angewandte Verfahren wurde 1835 durch ein anderes ersetzt, auf das Hofrat \textsc{Gauss} durch die Betrachtung der Induktionsgesetze gekommen war und das dem zuerst verwendeten bei Weitem vorzuziehen ist. Gerade bei dieser Art der Telegrafie hat der neue Apparat einen bedeutenden Vorzug vor dem gewöhnlichen Magnetometer. Das veranlasste Hofrat \textsc{Gauss}, dieses bis dahin öffentlich noch nicht erwähnte Verfahren im Vortrag in seinen Hauptzügen zu beschreiben und seine Leistungsfähigkeit anzugeben. Das müssen wir hier aus Platzgründen übergehen. Aus demselben Grund erwähnen wir auch nur kurz eine andere neue Vorrichtung, die jede unzeitige, bei einer bestimmten Arbeit störende Schwingung einer Magnetnadel in kurzer Zeit von selbst zur Ruhe bringen soll. Sie heißt deshalb Dämpfer. Diese Vorrichtung lässt sich beim neuen Apparat ebenso gut einsetzen wie beim Magnetometer und ist sowohl bei der erwähnten Art der Telegrafie als auch bei vielen anderen magnetischen Messarbeiten von wesentlichem Nutzen.

